\chapter{Oxygen therapy}
\index{Oxygen therapy}

\medskip
\noindent\textit{Educational information; use oxygen only as prescribed.}

\section*{What it is}
Oxygen therapy supplies a higher concentration of oxygen through nasal prongs, a mask, or another device. It is prescribed for documented low blood oxygen in selected conditions; it is not a general treatment for breathlessness and can be harmful when used without assessment.

\section*{Assessment and prescription}
A clinician decides whether oxygen is needed using symptoms, examination, oxygen measurements, and sometimes blood tests. The prescription specifies the flow or settings, when to use it, and the equipment. Do not change the flow, use someone else’s oxygen, or buy oxygen without medical advice.

\section*{Safe use at home}
Oxygen supports combustion. Do not smoke or allow flames, sparks, vaping, or heat sources near the equipment. Keep the room ventilated, store cylinders securely, keep tubing clear, and avoid oil-based products or flammable aerosols near the oxygen. Follow the supplier’s instructions and arrange equipment checks.

\section*{Possible problems}
Dryness, nose irritation, pressure injury, headache, or discomfort can occur. Too much oxygen can be dangerous for some people, particularly those prone to carbon-dioxide retention. Tell the care team about new sleepiness, confusion, worsening headache, or breathing changes.

\section*{When to seek urgent care}
Seek emergency help for severe or rapidly worsening breathlessness, blue or grey lips, confusion, collapse, or inability to use the prescribed equipment during a breathing emergency. Contact the prescribing team if the equipment fails, the tubing is damaged, or the person needs more oxygen than prescribed.

\section*{Sources}
\begin{itemize}
\item NHS home oxygen therapy: \url{https://www.nhs.uk/tests-and-treatments/home-oxygen-treatment/}
\end{itemize}
